% ===== §15 =====
\section{Making Room for the Other's Core}
\label{sec:room}

\noindent
This section draws from the persistence of drive a positive proposal about what a good bond consists in. Its aim is to establish that a relation is sustained not by the completeness of mutual understanding but by each partner's making room for the other's unsymbolizable core, and that the wish for total transparency, granted, would extinguish the very desire it sought to secure. The section states the proposal against the transparency ideal, grounds it in the thought experiment of Part Two, gives it as a claim, and turns the ordinary measure of intimacy around.

\subsection*{The proposal against the transparency ideal}

Each partner carries a core that was never symbolized and that goes on circling beneath their speech. The previous sections established that this core is the seat of the subject's desire, the void around which the drive turns, the very thing that makes them a desiring subject and not a transparent surface. A proposal about the good bond follows. A relation is not sustained by the success of communication, and its health is not measured by how completely the two have reached each other. This sets the account against the ideal, met already in the section on existence, that would make transparency the measure of a good bond, the thought that the best relation is the one in which nothing of either remains unsaid \citep{habermas1984theory}. That ideal, applied to intimacy, mistakes the good it aims at. To press toward the complete reaching of the other, the dissolution of every remainder, the full translation of the other into what one understands, is to press toward the extinction the thought experiment described, now enacted upon a person.

\subsection*{The ground in the thought experiment}

The ground of the proposal is the result of Part Two, transferred from the structure of the subject to the relation between two. The thought experiment showed that to close the gap is to extinguish the drive that circles it. An other wholly reached would be an other whose gap had been closed, whose remainder had been gathered into what the partner understands, and whose desire, therefore, had gone out. The wish to reach the other completely is thus, unknown to itself, a wish to still the other's desire, since the other's desire lives in the very core that complete reaching would dissolve. What presents itself as the deepest intimacy, a perfect mutual transparency with no remainder on either side, would be the death of desire on both, a bond of two surfaces that had been read and now had nothing further to turn on.

\subsection*{The good bond makes room}

The good bond does the opposite. It makes room for the other's core. It holds a space for what in the other cannot be said, does not press to close the gap, and lets the remainder go on circling as the living center of the other's desire. This making-room is not resignation, not a settling for less than understanding. It is the recognition that the unsymbolizable core is not a deficit in the bond but the presence in it of another desiring subject, and that to leave the core its room is to keep the other alive as other.

\begin{claim}[Room, not transparency]
A relation is sustained by each partner's making room for the other's unsymbolizable core, not by the completeness of mutual understanding. The core is the seat of the other's desire; to reach it wholly would be to close the gap the other's desire circles and so to extinguish it. What sustains a bond is two people each declining to fill what in the other must stay open, so that both remain subjects who desire rather than surfaces that have been read.
\end{claim}

\subsection*{The measure of intimacy turned around}

This turns the ordinary measure of intimacy around. The wish for complete mutual transparency, for a bond in which nothing of either remained unreached, is a wish that, granted, would end the desire that made the bond live. The deeper intimacy is the one that can dwell with the other's unsayable core without needing to close it, that finds in the very unreachability of the other not a barrier to love but the room in which the other stays a living source of it. What one loves in another, on this account, includes what one will never reach in them, and the not-reaching is not the frontier at which love stops but the space in which love continues. Part Four takes this making-room into three registers of relational life and shows the same generativity of the limit at work in the poem, in ethics, and in justice.
